\chapter{Аналитическая часть}

\section{Анализ предметной области}
Организация и проведение спортивных турниров -- это сложный многоуровневый процесс, требующий координации множества участников и ресурсов. Предметная область охватывает планирование соревнований, регистрацию участников, составление расписания матчей, назначение судей и фиксацию результатов.

Традиционно многие организаторы любительских и полупрофессиональных турниров используют для учета электронные таблицы (например, MS Excel) или бумажные носители. Это приводит к ряду проблем:
\begin{itemize}
	\item Высокая вероятность человеческой ошибки при вводе данных и подсчете очков;
	\item Сложность оперативного обновления турнирных таблиц и сеток;
	\item Отсутствие единого информационного пространства для игроков, судей и болельщиков;
	\item Трудности в оповещении участников об изменениях в расписании.
\end{itemize}

Внедрение автоматизированной системы управления турнирами (и соответствующей базы данных) позволит централизованно хранить информацию, автоматически создавать турнирные сетки, рассчитывать статистику и обеспечивать мгновенный доступ к результатам.

\subsection{Сравнительный анализ существующих конкурентов}
На российском рынке существует ряд популярных решений для управления спортивными турнирами. Для анализа выбраны три известные платформы: \textbf{Наградион}, \textbf{Goalstream} и \textbf{JoinSport}.

\begin{table}[htbp]
	\centering
	\caption{Сравнительный анализ систем управления турнирами}
	\vspace{0.2cm}
	\label{tab:competitors}
	\begin{tabular}{|p{3.5cm}|p{3.5cm}|p{3.5cm}|p{3.5cm}|}
		\hline
		\textbf{Критерий} & \textbf{Наградион} & \textbf{Goalstream} & \textbf{JoinSport} \\ \hline
		\textbf{Целевая аудитория} & Региональные федерации, крупные лиги. & Любительские командные виды спорта (чаще футбол). & Организаторы турниров, спортивные школы. \\ \hline
		\textbf{Поддержка форматов} & Максимально гибкая, для множества видов спорта. & Круговая система, плей-офф, группы. & Групповые этапы, плей-офф, шахматная система. \\ \hline
		\textbf{Пользовательские роли} & Администратор, судья, статистик, игрок. & Игрок, менеджер, организатор, болельщик. & Администратор, тренер, игрок, судья. \\ \hline
		\textbf{Интерфейс} & Перегруженный интерфейс из-за обилия настроек. & Удобный интерфейс с акцентом на социальные сети. & Современный интерфейс, конструктор сайтов. \\ \hline
		\textbf{Интеграция данных} & Выгрузка официальных протоколов, статистика. & Базовый экспорт, есть мобильное приложение. & Хороший API, виджеты для сторонних сайтов. \\ \hline
		\textbf{Недостатки} & Высокий порог вхождения, сложность для малых турниров. & Упор на функционал соцсети, узкая специализация. & Платный доступ к модулям, избыточность для разовых турниров. \\ \hline
	\end{tabular}
\end{table}


Анализ отечественных конкурентов показывает, что существующие решения либо слишком сложны и ориентированы на официальные федерации (Наградион), либо имеют узкую специализацию и перегружены социальными функциями (Goalstream). 
Проектируемая база данных и приложение должны взять лучшее из этих систем: гибкость в настройке форматов и четкое разделение ролей, при этом оставаясь интуитивно понятными для организаторов небольших и средних турниров. Основной упор в разрабатываемой БД будет сделан на оптимизацию структуры хранения данных команд, игроков, матчей и статистики, что обеспечит высокую скорость работы приложения без лишней функциональной перегруженности.

\section{Формализация задачи}

Необходимо спроектировать базу данных для хранения информации о пользователях, турнирах, командах, игроках, матчах, спортивных площадках и статистике. Требуется разработать приложение, предоставляющее интерфейс для просмотра, добавления, изменения и удаления информации, хранящейся в базе данных. Необходимо реализовать три вида ролей -- гость, представитель команды и организатор турниров.

%В рамках поставленной цели необходимо спроектировать и реализовать базу данных, содержащую информацию о турнирах, командах, расписании матчей и позволяющую изменять хранящуюся в ней информацию.


\section{Формализация данных}
Основываясь на анализе предметной области, можно выделить следующие категории данных:
\begin{itemize}
	\item пользователь;
	\item турнир;
	\item организатор турнира;
	\item команда;
	\item игрок;
	\item судья;
	\item тренер;
	\item матч;
	\item спортивная площадка.
\end{itemize}

Сведения о каждой категории данных содержится в таблице~\ref{tb:data}.

\begin{table}[ht]
	\begin{center}
		\begin{threeparttable}
			\caption{Категории и сведения о данных}
			\begin{tabular}{|c|p{10cm}|}
				\hline
				Категория & Сведения \\ \hline
				Организатор & Название компании, электронная почта, номер телефона \\ \hline
				Спортивная площадка & Название локации, город, вместимость \\ \hline
				Команда & Название команды, страна, дата основания \\ \hline
				Игрок & Привязка к команде, никнейм, настоящее имя, дата рождения \\ \hline
				Тренер & Привязка к команде, имя, фамилия, опыт работы (в годах) \\ \hline
				Судья & ФИО, категория, опыт работы (в годах) \\ \hline
				Турнир & Организатор, название, вид спорта, дата начала, дата окончания, призовой фонд, статус \\ \hline
				Заявка на турнир & Турнир, команда, дата регистрации, статус заявки (ожидает, принята, отклонена) \\ \hline
				Матч & Турнир, локация, дата и время проведения, стадия турнира, статус матча, команды-участницы и их счет, назначенные судьи и их роли \\ \hline
			\end{tabular}
			\label{tb:data} 
		\end{threeparttable}
	\end{center}
\end{table}
\clearpage

На рисунке~\ref{fig:er} приведена ER-диаграмма сущностей в нотации Чена.

\begin{figure}[h]
	\centering
	\includegraphics[width=0.9\textwidth]{img/er.png} % Не забудь добавить свою картинку по этому пути
	\caption{ER-диаграмма сущностей}
	\label{fig:er}
\end{figure}

Для корректного проведения турнира необходимо разработать систему учета матчей, которая основана на изменении статуса конкретной игры:
\begin{itemize}
	\item изначально при формировании расписания или турнирной сетки создается матч со статусом <<запланирован>>;
	\item по наступлению даты и времени начала матча, статус меняется на <<в игре>>;
	\item после внесения финального счета судьей или организатором матч переводится в статус <<завершен>>;
	\item в случае непредвиденных обстоятельств организатор может отменить игру, что приводит к изменению статуса на <<отменен>>;
\end{itemize}


\section{Формализация ролей}

В системе выделяются 3 типа пользователей: 
\begin{itemize}
	\item гость -- неавторизованный пользователь, обладающий возможностями просматривать список турниров, турнирные таблицы, расписание матчей, составы команд и статистику игроков;
	\item участник (представитель команды) -- авторизованный пользователь, обладающий возможностью просматривать общую информацию, подавать и редактировать заявку своей команды на участие в турнире, управлять списком игроков команды;
	\item организатор (администратор) --- авторизованный пользователь, обладающий возможностью просматривать, добавлять и менять данные турниров, команд, игроков, площадок, формировать расписание матчей, вносить результаты игр и управлять ролями пользователей.
\end{itemize}


На рисунке~\ref{fig:use-case} приведена Use-case диаграмма. 

\begin{figure}[h]
	\centering
	\includegraphics[width=1\textwidth]{usecase} 
	\caption{Use-case диаграмма}
	\label{fig:use-case}
\end{figure}


\section{Анализ моделей баз данных}

Основой любой современной информационной системы является база данных -- организованная структура для хранения, изменения и извлечения информации. Управление этой структурой осуществляется с помощью систем управления базами данных (СУБД), которые обеспечивают надежность, безопасность и многопользовательский доступ.

\subsection{Иерархическая и сетевая модели}
Исторически первыми моделями организации данных были иерархическая и сетевая. В иерархической модели данные представлены в виде строгой древовидной структуры, где каждая запись-потомок имеет строго одного родителя. Сетевая модель расширяет этот подход, позволяя записи иметь нескольких родителей, образуя сложный граф. 

Несмотря на высокую скорость навигации по жестко заданным физическим связям, эти подходы обладают существенным недостатком: они крайне негибки при изменении требований. В контексте спортивных турниров присутствуют сложные связи типа <<многие-ко-многим>> (например, команды участвуют в множестве турниров, а турниры включают множество команд; игроки могут менять клубы между сезонами). Описание таких взаимодействий с помощью иерархий сильно усложняет логику приложения, поэтому сегодня эти модели практически не применяются в новых проектах.

\subsection*{Нереляционные модели (NoSQL)}
В последние годы для решения специфических задач активно используются нереляционные модели данных. Наиболее популярным представителем является документоориентированная модель (например, СУБД MongoDB), где информация хранится в виде гибких JSON-документов. 

Документоориентированные БД отлично справляются с неструктурированной информацией и позволяют легко менять схему данных <<на лету>>. Однако для системы управления турнирами этот плюс оборачивается минусом. Отсутствие жесткой схемы повышает риск нарушения целостности данных (например, возможность сохранить матч без указания команд или с несуществующим турниром). Кроме того, документоориентированные базы хуже справляются со сложными транзакциями, затрагивающими множество сущностей одновременно.

\subsection*{Реляционная модель}
Реляционная модель, предложенная Э. Коддом, является стандартом для большинства современных учетных систем. Данные в ней представляются в виде двумерных таблиц (отношений), состоящих из строк (кортежей) и столбцов (атрибутов). Логические связи между таблицами устанавливаются с помощью первичных (Primary Key) и внешних (Foreign Key) ключей.

Ключевые преимущества реляционной модели для разрабатываемой системы:
\begin{itemize}
	\item \textbf{Строгая структурированность и целостность:} механизм внешних ключей не позволит удалить команду, если у нее есть сыгранные матчи в системе, что защищает историю турнира от разрушения.
	\item \textbf{Поддержка ACID-транзакций:} гарантирует, что процессы, состоящие из нескольких шагов (например, фиксация результатов матча, начисление очков командам и запись статистики игроков), выполняются атомарно -- либо полностью, либо никак.
	\item \textbf{Мощный аппарат агрегации (SQL):} реляционная алгебра идеально подходит для динамического расчета турнирного положения. С помощью SQL-запросов можно легко группировать данные, считать сумму забитых голов, разницу мячей и набранные очки.
\end{itemize}

\section*{Вывод по разделу}

В ходе анализа предметной области проведения спортивных соревнований были определены требования к разрабатываемой системе, выделены ключевые информационные сущности (турниры, команды, игроки, матчи) и формализованы роли пользователей.

Анализ специфики данных показал, что информация в спортивных турнирах обладает высокой степенью связности и строгой математической логикой. Матч не может существовать вне турнира, статистика привязана к конкретным игрокам, а результаты игр напрямую влияют на турнирное положение команд. 
Основываясь на необходимости обеспечения строгой ссылочной целостности, поддержки сложных аналитических агрегаций и транзакционной надежности, в качестве оптимальной модели хранения данных для разрабатываемой системы была выбрана \textbf{реляционная модель}.

\section*{Вывод}

<что сделали в аналите, кратко>

\clearpage
